\documentclass[12pt]{article}

\usepackage{bbm}
\usepackage{fancyhdr}
\usepackage{titling}
\usepackage{enumitem}
\usepackage{amsmath}
\usepackage{dsfont}
\usepackage{amsfonts}
\usepackage{amssymb}
\usepackage{amsthm}
\usepackage{geometry}
\usepackage{tikz}
\usepackage{algorithm}
\usepackage{algpseudocode}
\geometry{verbose,tmargin=1in,bmargin=1in,lmargin=1in,rmargin=1in}
\renewcommand{\vec}[1]{\mathbf{#1}}
\newcommand{\dec}{\text{dec}}
\newcommand{\divi}{\mathbin{\textnormal{div}}}
\title{Homework 1}
\author{Satvik Prasad}
\date{}
\newtheorem{lemma}{Lemma}[section]
\fancypagestyle{plain}{%
  \fancyhf{}
  \fancyhead[L]{\theauthor}
  \fancyhead[R]{\thedate}
}
\pagestyle{plain}
\makeatletter
\def\@maketitle{%
  \newpage
  \null
  \vskip 1em%
  \begin{center}%
    \let \footnote \thanks
    {\LARGE \@title \par}%
    \vskip 1em%
  \end{center}%
  \par
\vskip 1em}
\makeatother
\usepackage{cmbright}

\newcommand{\mat}[1]{\mathbf{#1}}

% ---------- Probability theory ----------
% Core operators (optional argument-free; use as \E{X}, \Var{X}, \Cov{X}{Y})
\DeclareMathOperator{\EOp}{\mathbb{E}}
\DeclareMathOperator{\VarOp}{Var}
\DeclareMathOperator{\CovOp}{Cov}
\DeclareMathOperator{\CorrOp}{Corr}

\renewcommand{\P}[1]{\mathbb{P}\left(#1\right)}        % \P{A}
\newcommand{\E}[1]{\EOp\left[#1\right]}                % \E{X}
\newcommand{\Var}[1]{\VarOp\left(#1\right)}            % \Var{X}
\newcommand{\Cov}[2]{\CovOp\left(#1,\,#2\right)}       % \Cov{X}{Y}
\newcommand{\Corr}[2]{\CorrOp\left(#1,\,#2\right)}     % \Corr{X}{Y}
\newcommand{\SD}[1]{\operatorname{SD}\left(#1\right)}  % \SD{X}

% Subscripted versions, e.g. \Esub{\theta}{X}, \Psub{\theta}{A}
\newcommand{\Esub}[2]{\EOp_{#1}\left[#2\right]}
\newcommand{\Psub}[2]{\mathbb{P}_{#1}\left(#2\right)}

% Conditioning, indicators, independence
\newcommand{\given}{\,\middle|\,}                      % \E{X \given Y}
\newcommand{\ind}[1]{\mathbbm{1}\left\{#1\right\}}     % \ind{X > 0}
\newcommand{\indep}{\perp\!\!\!\perp}                  % X \indep Y
\newcommand{\iid}{\overset{\text{iid}}{\sim}}          % X_i \iid \Normal(0,1)

% Convergence
\newcommand{\convp}{\xrightarrow{\;p\;}}               % in probability
\newcommand{\convd}{\xrightarrow{\;d\;}}               % in distribution
\newcommand{\convas}{\xrightarrow{\text{a.s.}}}        % almost surely
\newcommand{\convLp}[1]{\xrightarrow{L^{#1}}}          % in L^p

% Probability spaces
\newcommand{\F}{\mathcal{F}}                           % sigma-algebra
\newcommand{\B}{\mathcal{B}}                           % Borel sigma-algebra
\newcommand{\pspace}{(\Omega, \F, \mathbb{P})}

% Common distributions
\newcommand{\Normal}{\mathcal{N}}
\newcommand{\Bern}{\operatorname{Bernoulli}}
\newcommand{\Bin}{\operatorname{Binomial}}
\newcommand{\Pois}{\operatorname{Poisson}}
\newcommand{\Unif}{\operatorname{Uniform}}
\newcommand{\Exp}{\operatorname{Exp}}
\newcommand{\Geom}{\operatorname{Geometric}}
\newcommand{\Gam}{\operatorname{Gamma}}
\newcommand{\Beta}{\operatorname{Beta}}

% Number sets
\newcommand{\R}{\mathbb{R}}
\newcommand{\N}{\mathbb{N}}
\newcommand{\Z}{\mathbb{Z}}


\begin{document}
\begin{enumerate}
    \item First consider the case where $X$ is discrete.
    
    By the definition of the CDF, 
    $$
    F_X(s) = \sum_{k \leq s} P(X = k),
    $$
    and $$\lim_{\epsilon \downarrow 0}F_X(s - \epsilon) = \lim_{\epsilon \downarrow 0} \sum_{k \leq s - \epsilon} \mathbb P(X = k).$$

    Since for all $\epsilon > 0,$ $ s - \epsilon < s,$ $F_X(s)$ is not a part of the sum. Let $a$ be the largest real number satisfying $a < s$ and $a \in \text{range} \ X.$ Then, with $\epsilon = s - a,$ we have
    $$
        \sum_{k \leq s - \epsilon} \mathbb P(X = k) = \sum_{k \leq a} \mathbb P(X = k).
    $$
    For all $\epsilon < s- a,$ since $a$ is maximal, there is no $a < b < s$ such that $b \in \text{range}\ X,$ and thus $\sum_{k \leq s - \epsilon} \mathbb P(X = k) = \sum_{k \leq a} \mathbb P(X = k)$ for all $\epsilon < s - a.$ Therefore, by the definition of the limit, this limit $L$ is exactly defined by  
    $$
        L = \sum_{k \leq a } \mathbb P(X = k).
    $$

    Thus, $$F_X(s) - \lim_{e \downarrow } F_{X}(s -\epsilon) = \mathbb P(X = s)$$ for discrete $X.$

    Otherwise, we have $X$ continuous. Then, $\mathbb P(X = s) = 0,$ meaning we need to show that
    $$
        \lim_{\epsilon\downarrow 0} F_X(s - \epsilon ) = F_X(s).
    $$
    By the definition of a continuous random variable, $F_X$ is continuous, and thus left continuous, meaning this is trivially true. Hence, for continuous $X,$ this statement holds too.

    Finally, since $F_X(s)  = \lim_{\epsilon \downarrow 0}F_X(s + \epsilon),$ we have
    $$
        \mathbb P(X =s ) = F_X(s) - \lim_{\epsilon \downarrow 0} F_X(s -\epsilon) = \lim_{\epsilon \downarrow 0} F_X(s + \epsilon) + \lim_{\epsilon \downarrow 0} F_X(s - \epsilon)
    $$
    for all random variables $X$ as required. \qed

    \item We know $X$ has density function $f: \mathbb R \to \mathbb 
R$ defined by 
    $$
    f(x) = \dfrac{1}{2} e^{-|x|}.
    $$
    Then, $F_X(s)$ is defined by 
    $$
    F_X(s) = \dfrac{1}{2}\int_{-\infty}^{s} e^{-|x|} \, dx.
    $$

    When $x \leq 0,$ $f(x) = \frac{1}{2}e^{x},$ and $\frac{1}{2}e^{-x}$ otherwise. Hence, when $s \leq 0,$ we have
    $$
        F_X(s) = \dfrac{1}{2}\int_{-\infty}^{s} e^{x} \, dx = \dfrac{e^{s}}{2}.
    $$

    Otherwise, when $s > 0,$ we have
    \begin{align*}
    F_X(s) &= F_X(0) + \dfrac{1}{2}\int_{0}^{s} e^{-x} \, dx \\
    &= \dfrac{1}{2} + \dfrac{1 - e^{-s}}{2} = 1 - \dfrac{e^{-s}}{2}.
    \end{align*}

    Therefore, the distribution function is $F_X : \mathbb R \to [0, 1]$
    $$
        F_X(s) = \begin{cases}
            \dfrac{e^s}{2}  & s \leq 0 \\
            1 - \dfrac{e^{-s}}{2} & s > 0
        \end{cases} .
    $$
    \qed 
    
    \item We first find the distribution function $F_Y(s).$ We know
    $$
    \mathbb P(Y < s) = \mathbb P(|X| < s) = \mathbb P(-s < X < s),
    $$
    or equivalently, 
    $$
    \mathbb P(\{X < s\} \setminus \{X \leq -s\}) = F_X(s) - F_X(-s).
    $$
    Therefore, $f_Y(s)$ satisfies
    $$
        f_Y(s) = F'_Y(s) = f_x(s) + f_X(-s) = \dfrac{2}{\sqrt{2\pi}} e^{-s^2/2}.
    $$
    when $s \geq 0,$ and $0$ otherwise.

    Thus, the mean satisfies
    \begin{align*}
    \mu &= -\dfrac{2}{\sqrt{2\pi}}\int_{0}^{\infty}(-x)  e^{-x^2/2} \, dx \\
    &= - \dfrac{2}{\sqrt{2\pi}} [0 - 1] = \sqrt{\dfrac{2}{\pi}}
    \end{align*}

    For the variance, we compute
    $$
        E[Y^2] = \mathbb E[|X|^2] = \mathbb E[X^2] = \text{Var}[X] + \mathbb E[X]^2 = 1.
    $$
    Thus,
    $$
    \text{Var}[Y] = \mathbb E[Y^2] - \mu^2=  1- \dfrac{2}{\pi}.
    $$

    \item \begin{enumerate}
        \item Assume we have some $X \sim \text{Exp}(\lambda).$ Then, 
        $$
        F_X(s) = \begin{cases}
            1 - e^{-\lambda s} & s \geq 0, \\
            0 & \text{otherwise}
        \end{cases}.
        $$

        Then, for $s, t \geq 0,$ $$
        \mathbb P(X \geq s + t \mid X \geq s) = \dfrac{\mathbb P(X \geq s + t)}{\mathbb P(X \geq s)}.
        $$

        Hence, this is equal to 
        $$
            \dfrac{1 - F_X(s + t)}{1 - F_X(s)}  = \dfrac{e^{-\lambda(s + t)}}{e^{-\lambda(s)}} = e^{-\lambda t} = 1- (1 - e^{-\lambda t}) = \mathbb P(X \geq t)
        $$
        as required. \qed

        \item \begin{enumerate}
            \item Assume we have some continuous memoryless $X.$ Then, by definition, $$
                \mathbb P(X \geq s + t) = \mathbb P(X \geq t) \mathbb P(X \geq s).
            $$
            Since $mt = \underbrace{t + t + \ldots + t}_{\times m},$we can recurse $m$ times, yielding
            $$
            \mathbb P(X \geq mt) = \mathbb P(X \geq t)^{m}.
            $$

            \item Alternatively, consider $$
            \mathbb P(X \geq t/n)^{n} = \mathbb P(X \geq \underbrace{t/n + t/n + \ldots + t/n}_{t}) = \mathbb P(X \geq t),
            $$
            meaning
            $$
            \mathbb P(X \geq t/ n) = \mathbb P(X \geq t)^{1/n}.
            $$

            \item Then, assume we have some rational $q > 0.$ Hence, $q$ can be expressed as a ratio of $a, b > 0 \in \mathbb Z$ with $\gcd(a, b) = 1$. Thus, 
            $$
            \mathbb P(X \geq qt) = \mathbb P(X \geq at/b) = \mathbb P(X \geq at)^{1/b} = \mathbb P(X \geq t)^{a/b} = \mathbb P(X \geq t)^{q}
            $$
            as required.

            \item N.A.

            \item Thus, for every $s > 0,$ we can write
            $$
                \mathbb P(X \geq s) = \mathbb P(X \geq s(1)) = \mathbb P(X \geq 1)^{s}.
            $$
            Since $X$ is supported on $[0, \infty),$ $F_X(0) = 0.$ Using this result, we also know
            $$
            F_X(s) = 1- \mathbb P(X \geq s) = 1 - \mathbb P(X \geq 1)^{s}
            $$
            whenever $s > 0.$ Letting $\lambda = -\ln \mathbb P(X \geq 1),$ we get
            $$
            F_X(s) = \begin{cases}
                0 & s = 0 \\
                1 - e^{-\lambda s} & s > 0
            \end{cases} = 1-e^{-\lambda s}
            $$
            on $[0, \infty)$ as required. This is the distribution for $\text{Exp}(\lambda),$ meaning $X$ memoryless and continuous forces $X \sim \text{Exp}(\lambda)$ for some $\lambda = - \ln \mathbb P(X \geq 1).$
        \end{enumerate}
    \end{enumerate}

        \item \begin{enumerate}
        \item Let $X \sim \mathcal N(\mu, \sigma^2)$ and fix some $k > 0.$ Then, $X$ has density
        $$
        f_X(x) = \dfrac{1}{\sigma\sqrt{2\pi}} e^{-(x - \mu)^2/(2\sigma^2)},
        $$
        meaning
        $$
        \mathbb P(\mu - k\sigma \leq X \leq \mu + k\sigma) = \int_{\mu - k\sigma}^{\mu + k\sigma} \dfrac{1}{\sigma\sqrt{2\pi}} e^{-(x - \mu)^2/(2\sigma^2)} \, dx.
        $$
        Substituting $z = (x - \mu)/\sigma,$ we have $dx = \sigma \, dz,$ and the bounds $x = \mu \pm k\sigma$ become $z = \pm k.$ Hence,
        \begin{align*}
        \mathbb P(\mu - k\sigma \leq X \leq \mu + k\sigma) &= \int_{-k}^{k} \dfrac{1}{\sigma\sqrt{2\pi}} e^{-z^2/2} \, \sigma \, dz \\
        &= \int_{-k}^{k} \dfrac{1}{\sqrt{2\pi}} e^{-z^2/2} \, dz = \mathbb P(-k \leq \mathcal N(0, 1) \leq k),
        \end{align*}
        which does not depend on $\mu$ or $\sigma.$ Evaluating these integrals numerically for $k = 1, 2, 3$ gives
        \begin{align*}
        \mathbb P(-1 \leq \mathcal N(0, 1) \leq 1) &\approx 0.6827 \approx 68\%, \\
        \mathbb P(-2 \leq \mathcal N(0, 1) \leq 2) &\approx 0.9545 \approx 95\%, \\
        \mathbb P(-3 \leq \mathcal N(0, 1) \leq 3) &\approx 0.9973 \approx 99.7\%
        \end{align*}
        as required. \qed

        \item Assume $X \sim \text{Exp}(1).$ From the previous question, $\mathbb E[X] = 1/\lambda = 1$ and $\text{Var}[X] = 1/\lambda^2 = 1,$ so $\mu = \sigma = 1.$ Since $X$ is supported on $[0, \infty),$ and $\mu - k\sigma = 1 - k \leq 0$ for $k \geq 1,$ we have
        $$
        \mathbb P(\mu - k\sigma \leq X \leq \mu + k\sigma) = \mathbb P(0 \leq X \leq 1 + k) = \int_{0}^{1 + k} e^{-x} \, dx.
        $$
        Computing these numerically with
\begin{verbatim}
from scipy.integrate import quad
import math

lam = 1
mu = sigma = 1 / lam
for k in (1, 2, 3):
    p, _ = quad(lambda x: lam * math.exp(-lam * x),
                max(0, mu - k * sigma), mu + k * sigma)
    print(k, p)
\end{verbatim}
        gives
        \begin{align*}
        \mathbb P(0 \leq X \leq 2) &\approx 0.8647 \approx 86\%, \\
        \mathbb P(0 \leq X \leq 3) &\approx 0.9502 \approx 95\%, \\
        \mathbb P(0 \leq X \leq 4) &\approx 0.9817 \approx 98\%.
        \end{align*}

        \item Assume $X \sim \text{Exp}(2).$ Then, $\mu = \sigma = 1/2,$ and again $\mu - k\sigma \leq 0$ for $k \geq 1,$ so the intervals are $[0, 1], [0, 3/2], [0, 2].$ Setting \texttt{lam = 2} in the code above gives
        \begin{align*}
        \mathbb P(0 \leq X \leq 1) &= \int_0^1 2e^{-2x} \, dx \approx 0.8647 \approx 86\%, \\
        \mathbb P(0 \leq X \leq 3/2) &= \int_0^{3/2} 2e^{-2x} \, dx \approx 0.9502 \approx 95\%, \\
        \mathbb P(0 \leq X \leq 2) &= \int_0^2 2e^{-2x} \, dx \approx 0.9817 \approx 98\%,
        \end{align*}
        which are the same values as for $\text{Exp}(1).$

        \item There is a single rule. Assume $X \sim \text{Exp}(\lambda)$ for any $\lambda > 0.$ Then, $\mu = \sigma = 1/\lambda,$ and for $k \geq 1,$ $\mu - k\sigma = (1 - k)/\lambda \leq 0.$ Hence,
        \begin{align*}
        \mathbb P(\mu - k\sigma \leq X \leq \mu + k\sigma) &= \mathbb P\left(0 \leq X \leq \dfrac{1 + k}{\lambda}\right) \\
        &= F_X\left(\dfrac{1 + k}{\lambda}\right) = 1 - e^{-\lambda (1 + k)/\lambda} = 1 - e^{-(1 + k)},
        \end{align*}
        which does not depend on $\lambda.$

        For $k = 1, 2, 3,$ this gives
        $$
        1 - e^{-2} \approx 86.5\%, \quad 1 - e^{-3} \approx 95.0\%, \quad 1 - e^{-4} \approx 98.2\%,
        $$
        so it should be called something like the 86--95--98 rule.
    \end{enumerate}
\end{enumerate}
\end{document}